% New team-review material; existing manuscript sections are preserved.
\clearpage
\begingroup
\setlength{\tabcolsep}{5pt}
\setlength{\parskip}{0.35\baselineskip}
\section{High-school city evaluation: interim results}
\label{sec:t4progress}

Singapore is the conservative recommendation in our current comparison of
25 cities. The model separates a decision made with incomplete information
from a claim that one city is certainly best: all 25 cities remain possible
score leaders within the available evidence bounds. This section reports the
current stage of the city-evaluation study as of 8 November 2026.

\subsection{A decision model that retains missing information}
The eight main criteria describe safety, particulate air pollution, living
cost, housing affordability, upper secondary teaching resources, the supply
of higher education institutions, school travel time, and fixed broadband. Listed public
transport stops and public swimming pools form two auxiliary supply criteria.
Teaching resources measure student--teacher provision; they do not measure
learning outcomes. Institution and facility counts measure supply, with partial
inventories retained as lower bounds.

Let $u_{ij}\in[\ell_{ij},h_{ij}]\subseteq[0,1]$ denote city $i$'s utility
for criterion $j$. A completely unidentified value retains $[0,1]$.
For a nonnegative measured value $x$ and fixed reference $a>0$, benefit
criteria use $x/(a+x)$ and cost criteria use $a/(a+x)$. The reference values
remain fixed across cities. School travel instead uses
\begin{equation}
 u_T(t;\tau)=\exp(-t/\tau),\qquad \tau=30\ \text{minutes},
 \label{eq:t4travel}
\end{equation}
where $t$ is reported one-way duration. Time bins yield utility intervals;
unobserved journeys retain unknown utility. The scale $\tau$ expresses the
strength of the travel-time penalty and is tested below.

Each main criterion receives weight $(1-\alpha)/8$ and each auxiliary
criterion receives $\alpha/2$, with $\alpha=0.10$ for the reported ranking.
We use fixed ideal utilities of zero and one:
\begin{equation}
 C_i(u;\alpha)=100\,
 \frac{\sqrt{\sum_j w_j(\alpha)u_{ij}^{\,2}}}
 {\sqrt{\sum_j w_j(\alpha)u_{ij}^{\,2}}+
  \sqrt{\sum_j w_j(\alpha)(1-u_{ij})^2}}.
 \label{eq:t4score}
\end{equation}
The score increases with every utility, so its attainable bounds in the
rectangular evidence model are $L_i=C_i(\ell_i;\alpha)$ and
$U_i=C_i(h_i;\alpha)$. We rank choices by their worst-case regret:
\begin{equation}
 R_i=\max\{0,\max_{k\ne i}U_k-L_i\},\qquad
 i^*=\arg\min_i R_i.
 \label{eq:t4regret}
\end{equation}
Thus $R_i$ asks how far choosing city $i$ could leave a student behind a
better-scoring alternative. A city remains a possible score leader when
$U_i\geq\max_{k\ne i}L_k$. The rectangular model permits input values to vary
independently within their bounds; these are conditional decision statements.

\clearpage
\subsection{Current ranking and the evidence behind it}
Table~\ref{tab:t4rank} gives the decision order. Singapore's worst-case
regret is 34.07 score points, followed by Tokyo at 37.41. Phoenix ranks
16th. A numerical certificate retains Singapore as the minimum-regret
choice for every auxiliary weight $\alpha\in[0.05,0.20]$ under the current
inputs. The overlapping score intervals leave all 25 cities as possible
leaders; the decision order therefore expresses a conservative choice under
uncertainty.

\begin{table}[htbp]
\centering
\caption{Current minimum-regret ranking; smaller $R_i$ is preferred.
Scores use the evidence bounds and $\alpha=0.10$, $\tau=30$ minutes.}
\label{tab:t4rank}
\begin{tabular}{rlr@{\hspace{1.2em}}rlr}
\toprule
Rank & City & $R_i$ & Rank & City & $R_i$\\
\midrule
1 & Singapore & 34.07 & 14 & Delhi & 44.94 \\
2 & Tokyo & 37.41 & 15 & Istanbul & 45.09 \\
3 & Copenhagen & 38.10 & 16 & Phoenix & 46.14 \\
4 & Melbourne & 41.30 & 17 & Athens & 46.62 \\
5 & Cape Town & 42.02 & 18 & London & 46.78 \\
6 & Shanghai & 42.12 & 19 & Jakarta & 47.33 \\
7 & Sao Paulo & 42.67 & 20 & Nairobi & 48.24 \\
8 & Warsaw & 42.82 & 21 & New York & 48.42 \\
9 & Barcelona & 43.30 & 22 & Dhaka & 48.49 \\
10 & Toronto & 43.42 & 23 & Mumbai & 48.71 \\
11 & Bogota & 43.95 & 24 & Mexico City & 48.75 \\
12 & Dubai & 44.05 & 25 & Lagos & 51.58 \\
13 & Santiago & 44.67 &  & &  \\
\bottomrule
\end{tabular}
\end{table}

Actual student-travel information currently gives informative, conditional
intervals for five cities. New York uses a 2010--11 resident survey of
current grades 9--12 with an under-19 age restriction; unknown ages remain optional members. Mexico City uses
2025 private-household residents aged 14--18, with uncertain school-stage
membership retained. S\~ao Paulo uses 2023 municipal residents aged 14--18
in current upper-secondary education. Cape Town uses 2025 metropolitan
residents aged 14--18 in current high-school grades. Melbourne uses
2024--25 Greater Melbourne residents classified as Secondary, aged 15--19.
The other 20 cities retain $[0,1]$ for travel.

These sources differ in year, geography, age coverage and sampling procedure.
The S\~ao Paulo research mirror has not been verified as byte-identical to
the official release. Melbourne's reported interval is conditional on the
170 students with complete person weights; another eligible student's
unbounded missing weight leaves the strict full-group interval at $[0,1]$.
Source definitions, utility bounds and ranking inputs accompany this section
in the selected statistical snapshot.\footnote{Research model V108 and
validation V111; reproducible source records and calculations:
\url{https://github.com/TypeThe0ry/HiMCM/tree/1b35f6c/research/2026-a}.
The accompanying snapshot is \texttt{statistics/task4\_progress\_2026-11-08/}.}

\clearpage
\subsection{A revision that accounts for sampling uncertainty}
We added sampling uncertainty to Mexico City's 2025 time-bin interval,
holding the data and uncertain school-stage membership rules fixed. We retain
all 3,433 observed primary sampling units (PSUs) in 49 strata, including 586 units with
zero contribution to the target group. For each of nine fixed membership prefixes per endpoint, a Taylor-linearized
standard error produces a 97.5\% one-sided $t$ bound with 3,384 degrees of
freedom; taking
the minimum lower bound and maximum upper bound gives an approximate
95\% outer envelope for the identified interval. Final nonresponse-adjusted
weights are treated as fixed and finite-population corrections are omitted.

\begin{table}[htbp]
\centering
\caption{Mexico City before and after adding sampling uncertainty.
The other 249 city--criterion cells are unchanged.}
\label{tab:t4revision}
\begin{tabular}{lcc}
\toprule
Quantity & Before & After\\
\midrule
Travel utility interval & $[0.213824,0.584100]$ & $[0.207414,0.592051]$\\
Worst-case regret & 48.7028 & 48.7469\\
Decision rank & 23 & 24\\
\bottomrule
\end{tabular}
\end{table}

Mexico City moves down one place and Mumbai up one place; other ranks
remain unchanged. This corrects the treatment of uncertainty and leaves
25 possible leaders. School-stage interpretation and coverage remain
conditions of the estimate; the approximation has no exact finite-sample
coverage guarantee.

\subsection{Recommendation stability and rank sensitivity}
We re-evaluated the five informative travel sources at five prespecified
scales, retaining their original membership, missing-data and sampling rules.
Each scale reproduces utilities from the available duration or time-bin
information. Kendall's $\tau_b$ compares each decision order with the
default-scale order.

\begin{table}[htbp]
\centering
\caption{Direct travel-scale sensitivity relative to $\tau=30$ minutes.}
\label{tab:t4sensitivity}
\begin{tabular}{lrrrrr}
\toprule
$\tau$ (minutes) & 15 & 20 & 30 & 45 & 60\\
\midrule
Cities changing rank & 11 & 9 & 0 & 9 & 10\\
Kendall's $\tau_b$ & 0.9267 & 0.9533 & 1.0000 & 0.9600 & 0.9400\\
\bottomrule
\end{tabular}
\end{table}

Singapore remains the conservative recommendation across all five scales
and the continuous $\alpha$ range at each scale. Cape Town ranks fourth
to ninth and Mexico City nineteenth to twenty-fourth. All 25 cities remain
possible leaders. Thus the recommendation is stable across the five tested
scales while intermediate ranks vary. Values between the tested $\tau$
settings were not continuously certified.
\clearpage
\endgroup
